%----------
%	ANÁLISIS EXPERTO Y PERSPECTIVA DE GÉNERO — proyecto IRIS
%----------

El presente trabajo se enmarca en el proyecto IRIS del Instituto Universitario
de Estudios de Género de la Universidad Carlos~III de Madrid, cuyo objetivo es
analizar de forma sistemática la representación de las mujeres en la prensa
española. A diferencia de aproximaciones que abordan la detección de discurso
de odio sobre mensajes breves en redes sociales, el objeto de estudio aquí son
\textbf{piezas periodísticas completas} y el fenómeno analizado no es la
agresión explícita, sino el \textbf{sexismo lingüístico}: los usos del lenguaje
que reproducen desigualdad de forma frecuentemente inadvertida por quien los
emplea.

\subsection{Fundamentación metodológica}

El libro de códigos empleado se apoya en la tradición del análisis de contenido
con perspectiva de género aplicada a medios de comunicación
\cite{sainzdebaranda2013, sainzdebaranda2026}, y operacionaliza criterios
procedentes de guías institucionales de lenguaje no sexista elaboradas por
organismos públicos españoles. Estas guías proporcionan tanto los principios
generales como la casuística concreta a partir de la cual se han derivado las
reglas de decisión de cada variable.

Un concepto resulta central en esta fundamentación: la \textbf{regla de
inversión}. Ante un pasaje candidato, el procedimiento consiste en reescribirlo
intercambiando el género de todas las referencias personales y evaluar si el
resultado se publicaría con naturalidad. Si la versión invertida resulta
extraña, ridícula o irrespetuosa, cabe concluir que la formulación original
incorpora sexismo lingüístico. Este criterio permite distinguir con precisión
entre textos que \emph{hablan sobre} desigualdad —informativos y no
necesariamente sexistas— y textos que \emph{reproducen} la desigualdad en su
forma.

Conviene subrayar que el sexismo lingüístico opera sobre la forma del mensaje y
no sobre su fondo. Una pieza puede presentar un contenido no sexista expresado
mediante formas sexistas, o a la inversa. Esta distinción justifica que el libro
de códigos separe explícitamente las variables de forma de aquellas de
contenido, tal como se detalla en la Sección~\ref{sec:iris_variables}.

\subsection{Procedimiento de anotación}

La anotación del corpus fue realizada por especialistas en comunicación y
estudios de género vinculadas al proyecto, aplicando el libro de códigos sobre
la pieza completa. La unidad de análisis es, por tanto, \textbf{el artículo}: cada
variable recibe un único veredicto por pieza, con independencia del número de
ocurrencias del fenómeno en el texto. Esta decisión, coherente con la práctica
habitual en análisis de contenido, tiene implicaciones relevantes para la
interpretación de los resultados, pues basta una sola ocurrencia para que una
pieza extensa quede codificada como positiva.

\subsection{Limitación: ausencia de doble codificación}
\label{subsec:limitacion_fiabilidad}

Cada noticia del corpus fue codificada por una única anotadora. En consecuencia,
\textbf{no es posible calcular la fiabilidad intercodificadora} del libro de
códigos sobre este material, ni estimar por tanto el techo de acuerdo alcanzable
por un sistema automático.

Esta limitación debe tenerse presente en la interpretación de los resultados del
Capítulo~\ref{results_and_discussion}. Las variables consideradas implican
juicios interpretativos sobre los que dos personas expertas podrían discrepar
legítimamente; en tareas de análisis de contenido con este grado de
subjetividad, los coeficientes de acuerdo entre codificadoras humanas raramente
alcanzan valores próximos a la unidad. Al carecer de esa referencia, cualquier
discrepancia observada entre el sistema y la anotación experta \emph{no puede
atribuirse inequívocamente al error del modelo}: parte de ella podría
corresponder a la variabilidad inherente a la propia tarea de codificación.

Se recomienda, como línea de trabajo futura de bajo coste y alto valor, la doble
codificación de una submuestra del corpus (habitualmente entre el 10\,\% y el
20\,\% en la literatura de análisis de contenido), suficiente para obtener una
estimación defendible de dicha fiabilidad.
